\documentclass[11pt,a4paper]{article}

\usepackage[T1]{fontenc}
\usepackage[utf8]{inputenc}
\usepackage[indonesian]{babel}
\usepackage{amsmath,amssymb}
\usepackage{booktabs,longtable,array}
\usepackage{enumitem}
\usepackage[hidelinks]{hyperref}
\usepackage[margin=2.4cm]{geometry}

\hypersetup{
  pdftitle={O013: Rute Keterhubungan dan Penguasaan},
  pdfauthor={OpenAI Codex gpt-5.6-sol, Ultra},
  pdfsubject={Lapisan orisinal penghubung dan penguasaan untuk edisi Bahasa Indonesia},
  pdfkeywords={aljabar pascasarjana, teori representasi, aljabar komutatif, Bahasa Indonesia}
}

\setlist{nosep}
\newcommand{\oid}[1]{\texttt{#1}}
\newcommand{\jalur}[1]{\textbf{#1}}

\title{O013: Rute Keterhubungan dan Penguasaan\\
\large Dari Metode Aljabar ke Teori Representasi dan Aljabar Komutatif}
\author{OpenAI Codex gpt-5.6-sol, Ultra\\\small disusun atas instruksi pengguna}
\date{29 Agustus 2026}

\begin{document}
\maketitle

\section*{Status orisinal dan provenance}

\textbf{Status:} seluruh unit O013 ini merupakan materi orisinal edisi.
Unit ini bukan terjemahan, ringkasan, parafrasa, atau bagian dari akar sumber
mana pun. Nama komponen pembelajaran di bawah hanya berfungsi sebagai tujuan
navigasi. Tidak ada atribusi kepengarangan sumber yang diterapkan pada materi
orisinal ini.

\textbf{Provenance produksi:} \emph{OpenAI Codex gpt-5.6-sol, Ultra, on instructions of the user}.
Unit disusun khusus sebagai lapisan penghubung dan
penguasaan bagi edisi bahasa Indonesia ini. Backend ber-ID stabil terdapat
dalam \texttt{o013-rute-pembelajar.json}.

\section{Tujuan dan cara menggunakan unit}

Unit ini membantu pembelajar berpindah dari landasan aljabar umum menuju dua
jalur yang saling melengkapi:

\begin{enumerate}
\item \jalur{jalur R}, yakni teori representasi grup melalui representasi,
modul, teori karakter, induksi, grup simetrik, dan medan dasar yang tidak
tertutup; serta
\item \jalur{jalur C}, yakni enam topik aljabar komutatif terpilih: Lema
Nakayama; spektrum dan topologi Zariski; prima terkait dan dekomposisi primer;
teorema terletak di atas dan naik; Nullstellensatz dan normalisasi; serta
dimensi Krull.
\end{enumerate}

Rute ini bukan perintah untuk membaca semuanya secara linear. Kerjakan
diagnostik terlebih dahulu, perbaiki prasyarat yang gagal, lalu ikuti simpul
dependensi. Suatu simpul dianggap dikuasai apabila pembelajar dapat menjelaskan
definisi utamanya, menghasilkan satu contoh dan satu bukan-contoh, serta
menyelesaikan prompt penguasaan yang terkait tanpa membuka jawaban.

\section{Prasyarat yang perlu tersedia}

\begin{longtable}{>{\raggedright\arraybackslash}p{2.3cm}p{5.0cm}p{7.0cm}}
\toprule
ID & Prasyarat & Bukti kesiapan minimum \\
\midrule
\endhead
\oid{O013-P01} & Bahasa fungsi, hasil bagi, dan pembuktian & Dapat membedakan
definisi, implikasi, dan konvers; dapat memeriksa bahwa suatu pemetaan pada
kelas ekuivalensi terdefinisi dengan baik. \\
\oid{O013-P02} & Grup, subgrup, hasil bagi, dan aksi grup & Dapat menghitung
orbit dan penstabil, menggunakan teorema orbit--penstabil, serta bekerja dengan
koset dan konjugasi. \\
\oid{O013-P03} & Gelanggang komutatif, ideal, dan pelokalan & Dapat membedakan
ideal prima dan maksimal, menghitung ideal hasil bagi sederhana, dan
menjelaskan apa yang dibalik dalam $S^{-1}R$. \\
\oid{O013-P04} & Modul, ruang vektor, tensor, dan barisan eksak & Dapat
menggunakan hasil kali tensor, mengenali modul terbangkitkan hingga, membaca
barisan eksak pendek, dan menjelaskan jumlah langsung serta modul
semisederhana. \\
\oid{O013-P05} & Aljabar, integralitas, norma, dan teras & Dapat menguji bahwa
suatu unsur integral dengan persamaan monik dan membedakan perluasan hingga,
aljabar, dan transenden. \\
\oid{O013-P06} & Ekstensi medan dan penutup aljabar & Dapat bekerja dengan
polinomial minimal, medan pemisah, dan makna medan tertutup secara aljabar. \\
\bottomrule
\end{longtable}

Prasyarat \oid{O013-P01} bersifat lintas jalur. Jalur R terutama memakai
\oid{O013-P02} dan \oid{O013-P04}; jalur C terutama memakai
\oid{O013-P03}--\oid{O013-P06}. Pembelajar tidak perlu menuntaskan jalur R
sebelum memulai jalur C, atau sebaliknya.

\section{Peta dependensi eksplisit}

\subsection*{Simpul jalur R: teori representasi}

\begin{longtable}{>{\raggedright\arraybackslash}p{2.3cm}p{5.2cm}p{6.8cm}}
\toprule
ID & Sasaran pembelajaran & Memerlukan \\
\midrule
\endhead
\oid{O013-R10} & Representasi linear dan aljabar grup &
\oid{O013-P02}, \oid{O013-P04} \\
\oid{O013-R11} & Aljabar multilinear, tensor, dual, simetrik, dan eksterior &
\oid{O013-P04}, \oid{O013-P05} \\
\oid{O013-R12} & Modul, semisederhana, dan teori Wedderburn &
\oid{O013-P03}, \oid{O013-P04}, \oid{O013-R10} \\
\oid{O013-R13} & Gelanggang representasi, karakter, dan ortogonalitas &
\oid{O013-R10}, \oid{O013-R11}, \oid{O013-R12} \\
\oid{O013-R14} & Restriksi, representasi terinduksi, dan resiprositas Frobenius &
\oid{O013-P02}, \oid{O013-R13} \\
\oid{O013-R15} & Grup simetrik dan grup linear umum &
\oid{O013-R11}, \oid{O013-R14} \\
\oid{O013-R16} & Representasi atas medan yang tidak tertutup &
\oid{O013-P06}, \oid{O013-R12}, \oid{O013-R13} \\
\bottomrule
\end{longtable}

Tulang punggung jalur ini adalah
\[
\oid{O013-P02}+\oid{O013-P04}
\longrightarrow \oid{O013-R10}
\longrightarrow \oid{O013-R12}
\longrightarrow \oid{O013-R13}
\longrightarrow \oid{O013-R14}.
\]
Simpul \oid{O013-R11} memasok alat multilinear bagi
\oid{O013-R13} dan \oid{O013-R15}; simpul \oid{O013-R16} merupakan cabang
yang memerlukan kesiapan medan.

\subsection*{Simpul jalur C: enam topik aljabar komutatif}

\begin{longtable}{>{\raggedright\arraybackslash}p{2.3cm}p{5.2cm}p{6.8cm}}
\toprule
ID & Sasaran pembelajaran & Memerlukan \\
\midrule
\endhead
\oid{O013-C10} & Lema Nakayama dan pembangkit minimal modul lokal &
\oid{O013-P03}, \oid{O013-P04} \\
\oid{O013-C11} & $\operatorname{Spec}(R)$ dan topologi Zariski &
\oid{O013-P03} \\
\oid{O013-C12} & Prima terkait, modul primer, dan dekomposisi primer &
\oid{O013-C10}, \oid{O013-C11}, \oid{O013-P04} \\
\oid{O013-C13} & Teorema terletak di atas dan naik untuk perluasan integral &
\oid{O013-P05}, \oid{O013-P06}, \oid{O013-C11} \\
\oid{O013-C14} & Nullstellensatz Hilbert dan normalisasi &
\oid{O013-P05}, \oid{O013-P06}, \oid{O013-C11}, \oid{O013-C13} \\
\oid{O013-C15} & Fungsi Hilbert dan dimensi Krull &
\oid{O013-C10}, \oid{O013-C11}, \oid{O013-C12}, \oid{O013-C14} \\
\bottomrule
\end{longtable}

Tulang punggung jalur ini adalah
\[
\oid{O013-P03}\longrightarrow\oid{O013-C11},\qquad
\oid{O013-P03}+\oid{O013-P04}\longrightarrow\oid{O013-C10},
\]
diikuti dua cabang:
\[
\oid{O013-C10}+\oid{O013-C11}\longrightarrow\oid{O013-C12}
\]
dan
\[
\oid{O013-P05}+\oid{O013-P06}+\oid{O013-C11}
\longrightarrow\oid{O013-C13}
\longrightarrow\oid{O013-C14}.
\]
Kedua cabang bertemu kembali pada \oid{O013-C15}.

\subsection*{Simpul sintesis}

\oid{O013-S10} menghubungkan kedua jalur setelah
\oid{O013-R12} dan \oid{O013-C11}. Gagasan penghubungnya adalah bahwa suatu
representasi dapat dipandang sebagai modul atas aljabar grup, sedangkan
anihilator, dukungan, dan pelokalan menjelaskan bagaimana suatu modul terlihat
dari berbagai ideal prima. Hubungan ini bersifat struktural; karakter grup dan
titik spektrum bukan objek yang sama dan tidak boleh disamakan.

\section{Urutan studi yang disarankan}

\begin{enumerate}[label=\textbf{Tahap \arabic*.}]
\item \textbf{Masuk dan memilah.} Kerjakan \oid{O013-D01}--\oid{O013-D08}.
Catat bukan hanya skor, tetapi juga alasan setiap jawaban.
\item \textbf{Perbaikan landasan.} Ulangi hanya prasyarat yang ditunjukkan
oleh diagnostik. Jangan menunda seluruh rute karena satu subtopik kecil.
\item \textbf{Tulang punggung representasi.} Baca dan kerjakan
\oid{O013-R10}, \oid{O013-R12}, \oid{O013-R13}, lalu
\oid{O013-R14}. Sisipkan \oid{O013-R11} sebelum perhitungan karakter yang
memakai tensor, dual, atau pangkat simetrik.
\item \textbf{Cabang representasi.} Ambil \oid{O013-R15} untuk fokus
kombinatorial atau \oid{O013-R16} untuk fokus pada pengaruh medan dasar.
Keduanya boleh dipelajari dalam urutan mana pun setelah dependensinya terpenuhi.
\item \textbf{Tulang punggung komutatif.} Kerjakan \oid{O013-C10} dan
\oid{O013-C11}; keduanya dapat berjalan paralel setelah prasyaratnya siap.
\item \textbf{Struktur dan geometri.} Lanjutkan ke \oid{O013-C12}, serta
secara terpisah ke \oid{O013-C13} lalu \oid{O013-C14}.
\item \textbf{Dimensi dan sintesis.} Kerjakan \oid{O013-C15}, kemudian
\oid{O013-S10} dan prompt penguasaan yang belum tuntas.
\end{enumerate}

Untuk setiap simpul, gunakan siklus empat langkah: baca definisi; buat contoh;
rekonstruksi satu pembuktian tanpa melihat; lalu kerjakan satu prompt. Jika
gagal pada langkah ketiga, kembali ke dependensi langsung, bukan otomatis ke
awal buku.

\section{Diagnostik masuk}

Jawablah tanpa catatan. Beri satu poin hanya jika kesimpulan dan alasan sama-sama
benar.

\begin{description}[leftmargin=3.2cm,style=nextline]
\item[\oid{O013-D01}] Untuk aksi $G$ pada $X$ dan $x\in X$, jelaskan hubungan antara
$|Gx|$, $|G|$, dan $|\operatorname{Stab}_G(x)|$. Kapan rumus kardinalitas ini
berlaku?
\item[\oid{O013-D02}] Misalkan $k$ medan dan $G$ grup hingga. Tuliskan perkalian dua unsur
umum $\sum_g a_g g$ dan $\sum_h b_h h$ dalam aljabar grup $kG$.
\item[\oid{O013-D03}] Nyatakan sifat universal $V\otimes_k W$ dalam satu kalimat dan jelaskan
mengapa pemetaan bilinear $V\times W\to U$ menghasilkan pemetaan linear dari
$V\otimes_k W$.
\item[\oid{O013-D04}] Jika $M$ modul terbangkitkan hingga atas gelanggang lokal $(R,\mathfrak
m)$, informasi apa tentang pembangkit $M$ yang tersimpan dalam ruang vektor
$M/\mathfrak mM$?
\item[\oid{O013-D05}] Bedakan ideal prima, ideal maksimal, dan ideal radikal. Berikan satu
implikasi yang selalu benar dan satu konvers yang dapat gagal.
\item[\oid{O013-D06}] Apa arti suatu unsur $s$ integral atas subgelanggang $R$? Mengapa syarat
``monik'' penting?
\item[\oid{O013-D07}] Jika $T$ operator linear berorde hingga pada ruang vektor kompleks,
mengapa $T$ dapat didiagonalkan dan apa bentuk nilai eigennya?
\item[\oid{O013-D08}] Berikan definisi panjang rantai ideal prima dan jelaskan bagaimana rantai
tersebut memberi batas bawah bagi dimensi Krull.
\end{description}

\paragraph{Indikator jawaban.}
\oid{O013-D01}: teorema orbit--penstabil untuk grup hingga.
\oid{O013-D02}: koefisien pada $u\in G$ adalah jumlah $a_gb_h$ atas
$gh=u$. \oid{O013-D03}: faktorisasi tunggal pemetaan bilinear melalui pemetaan
kanonik ke tensor. \oid{O013-D04}: dimensi $M/\mathfrak mM$ adalah banyaknya
pembangkit minimal. \oid{O013-D05}: maksimal mengimplikasikan prima dan prima
mengimplikasikan radikal, tetapi konversnya dapat gagal. \oid{O013-D06}:
persamaan polinomial monik memberi kontrol keterhinggaan. \oid{O013-D07}:
polinomial minimal membagi $X^n-1$, yang berakar sederhana dalam
$\mathbb C$. \oid{O013-D08}: panjang adalah banyaknya inklusi ketat; supremumnya
memberi dimensi.

\paragraph{Keputusan masuk.}
Skor $0$--$3$: perbaiki \oid{O013-P01}--\oid{O013-P04} sebelum memilih
cabang. Skor $4$--$6$: mulai rute, tetapi pasangkan setiap jawaban gagal dengan
prasyaratnya. Skor $7$--$8$: masuk langsung. Terlepas dari skor total,
\oid{O013-D01}--\oid{O013-D03} adalah gerbang jalur R, sedangkan
\oid{O013-D04}--\oid{O013-D06} adalah gerbang jalur C.

\section{Prompt penguasaan orisinal}

Semua prompt, petunjuk, dan jawaban pada bagian ini baru ditulis untuk O013.
Gunakan petunjuk secara berurutan dan buka jawaban hanya setelah menuliskan
upaya lengkap.

\subsection*{\oid{O013-M01}: Representasi sebagai modul aljabar grup}

\textbf{Prompt.} Misalkan $k$ medan dan $G$ grup hingga. Buktikan bahwa memberi
representasi $\rho:G\to\operatorname{GL}(V)$ ekuivalen dengan memberi struktur
modul kiri $kG$ pada $V$. Tunjukkan pula bahwa morfisme pada kedua sisi
berkorespondensi.

\textbf{Petunjuk \oid{O013-M01-H1}.} Perluas aksi setiap $g\in G$ secara
$k$-linear ke kombinasi formal unsur grup.

\textbf{Petunjuk \oid{O013-M01-H2}.} Untuk arah balik, batasi aksi $kG$ ke
basis grup $G\subset kG$ dan gunakan $g^{-1}$ untuk menunjukkan bahwa aksinya
dapat dibalik.

\textbf{Jawaban \oid{O013-M01-A}.} Definisikan
\[
\left(\sum_{g\in G}a_gg\right)v=\sum_{g\in G}a_g\rho(g)v.
\]
Sifat modul mengikuti dari linearitas dan $\rho(gh)=\rho(g)\rho(h)$.
Sebaliknya, struktur modul memberi $\rho(g)v=gv$; relasi dalam $kG$ memberi
$\rho(gh)=\rho(g)\rho(h)$, dan $g^{-1}$ memberi invers. Pemetaan linear
$f:V\to W$ adalah $G$-ekuivarian tepat ketika $f(gv)=gf(v)$ untuk setiap
$g$, dan melalui linearitas hal ini tepat berarti $f$ linear terhadap $kG$.

\subsection*{\oid{O013-M02}: Proyeksi perataan dan ketereduksian lengkap}

\textbf{Prompt.} Misalkan $\operatorname{char}(k)$ tidak membagi $|G|$,
$V$ representasi berdimensi hingga, dan $W\subseteq V$ subrepresentasi.
Mulailah dengan sebarang proyeksi linear $p:V\to W$. Bentuk proyeksi
$G$-ekuivarian dan deduksikan $V=W\oplus W'$ untuk suatu subrepresentasi $W'$.

\textbf{Petunjuk \oid{O013-M02-H1}.} Rata-ratakan semua konjugat dari $p$:
$gpg^{-1}$.

\textbf{Petunjuk \oid{O013-M02-H2}.} Periksa secara terpisah bahwa hasil
perataan tetap identitas pada $W$ dan komutatif dengan aksi setiap $h\in G$.

\textbf{Jawaban \oid{O013-M02-A}.} Ambil
\[
\overline p=\frac1{|G|}\sum_{g\in G}gpg^{-1}.
\]
Penyebut sah menurut asumsi karakteristik. Pengindeksan ulang $g\mapsto h^{-1}g$
menunjukkan $h\overline p=\overline ph$. Karena $W$ stabil dan $p|_W$ identitas,
setiap $gpg^{-1}$ juga identitas pada $W$; jadi $\overline p$ adalah proyeksi
ke $W$. Kernel $W'=\ker\overline p$ stabil terhadap $G$, sehingga
$V=W\oplus W'$.

\subsection*{\oid{O013-M03}: Membaca multiplisitas dari karakter}

\textbf{Prompt.} Untuk $S_3$, gunakan urutan kelas identitas, transposisi, dan
siklus-$3$, dengan kardinalitas kelas $1,3,2$. Karakter tak tereduksi adalah
\[
\mathbf{1}=(1,1,1),\qquad \varepsilon=(1,-1,1),\qquad \sigma=(2,0,-1).
\]
Suatu representasi memiliki karakter $\chi=(5,1,-1)$. Tentukan dekomposisinya.

\textbf{Petunjuk \oid{O013-M03-H1}.} Gunakan hasil kali dalam berbobot
$\langle\alpha,\beta\rangle=\frac16\sum_C|C|\overline{\alpha(C)}\beta(C)$.

\textbf{Petunjuk \oid{O013-M03-H2}.} Hitung tiga hasil kali dalam dan periksa
dimensi pada unsur identitas.

\textbf{Jawaban \oid{O013-M03-A}.} Diperoleh
\[
\langle\mathbf1,\chi\rangle=1,\qquad
\langle\varepsilon,\chi\rangle=0,\qquad
\langle\sigma,\chi\rangle=2.
\]
Jadi $V\cong\mathbf1\oplus\sigma^{\oplus2}$. Pemeriksaan dimensi memberi
$1+2\cdot2=5=\chi(1)$.

\subsection*{\oid{O013-M04}: Nakayama dan pembangkit minimal}

\textbf{Prompt.} Misalkan $(R,\mathfrak m,k)$ gelanggang lokal dan $M$ modul
terbangkitkan hingga. Jika kelas $\overline m_1,\ldots,\overline m_r$ membentuk
basis-$k$ dari $M/\mathfrak mM$, buktikan bahwa $m_1,\ldots,m_r$ merupakan
sistem pembangkit minimal dari $M$.

\textbf{Petunjuk \oid{O013-M04-H1}.} Bentuk $N=M/(Rm_1+\cdots+Rm_r)$ dan
hitung $N/\mathfrak mN$.

\textbf{Petunjuk \oid{O013-M04-H2}.} Terapkan Lema Nakayama pada $N$, lalu
gunakan dimensi hasil bagi untuk membuktikan minimalitas.

\textbf{Jawaban \oid{O013-M04-A}.} Kelas-kelas tersebut merentang hasil bagi,
sehingga $N/\mathfrak mN=0$, yakni $N=\mathfrak mN$. Modul $N$ terbangkitkan
hingga; Lema Nakayama memberi $N=0$. Jadi semua $m_i$ membangkitkan $M$.
Setiap sistem pembangkit $s$ unsur memetakan ke himpunan perentang
$M/\mathfrak mM$, sehingga $s\ge r$. Sistem yang ditemukan minimal.

\subsection*{\oid{O013-M05}: Spektrum dan pelokalan}

\textbf{Prompt.} Untuk $f\in R$, buktikan bahwa kontraksi memberi bijeksi
\[
\operatorname{Spec}(R_f)\longleftrightarrow
D(f)=\{\mathfrak p\in\operatorname{Spec}(R):f\notin\mathfrak p\}.
\]
Uji hasilnya untuk $R=k[x,y]$, $f=x$, dengan ideal $(y)$ dan $(x,y)$.

\textbf{Petunjuk \oid{O013-M05-H1}.} Gunakan korespondensi ideal prima pada
$S^{-1}R$ dengan ideal prima pada $R$ yang tidak bertemu $S=\{1,f,f^2,\ldots\}$.

\textbf{Petunjuk \oid{O013-M05-H2}.} Periksa ekstensi lalu kontraksi, serta
kontraksi lalu ekstensi.

\textbf{Jawaban \oid{O013-M05-A}.} Jika $\mathfrak q$ prima dalam $R_f$,
kontraksinya tidak memuat $f$ karena $f$ satuan. Jika $\mathfrak p$ prima dan
$f\notin\mathfrak p$, maka $\mathfrak pR_f$ prima dan kedua operasi saling
invers. Untuk contoh, $(y)$ berada dalam $D(x)$ dan memberi ideal prima
$(y)R_x$, sedangkan $(x,y)$ memuat $x$ sehingga tidak memberi titik
$\operatorname{Spec}(R_x)$.

\subsection*{\oid{O013-M06}: Prima terkait dan dekomposisi primer}

\textbf{Prompt.} Dalam $R=k[x,y]$, ambil $I=(x^2,xy)$. Verifikasi
\[
I=(x)\cap(x^2,y),
\]
tentukan radikal kedua komponen, dan simpulkan kandidat untuk
$\operatorname{Ass}_R(R/I)$.

\textbf{Petunjuk \oid{O013-M06-H1}.} Untuk inklusi yang sukar, tulis unsur
irisan sebagai $ax^2+by$ dan gunakan keterbagian oleh $x$.

\textbf{Petunjuk \oid{O013-M06-H2}.} Ideal $(x)$ prima; hasil bagi oleh
$(x^2,y)$ isomorfik dengan $k[x]/(x^2)$.

\textbf{Jawaban \oid{O013-M06-A}.} Inklusi $I$ ke irisan langsung. Jika
$ax^2+by$ berada dalam $(x)$, maka $x\mid by$; karena $x$ dan $y$ relatif
prima, $x\mid b$, sehingga unsur tersebut berada dalam $(x^2,xy)$. Jadi
persamaan benar. Komponen $(x)$ adalah $(x)$-primer, sedangkan $(x^2,y)$
adalah $(x,y)$-primer dengan radikal $(x,y)$. Dekomposisi ini memberi
\[
\operatorname{Ass}_R(R/I)=\{(x),(x,y)\}.
\]

\subsection*{\oid{O013-M07}: Integralitas, normalisasi, dan dimensi}

\textbf{Prompt.} Misalkan $A=k[t^2,t^3]\subset B=k[t]$. Buktikan bahwa
$B$ integral dan terbangkitkan hingga atas $A$, jelaskan surjektivitas
$\operatorname{Spec}(B)\to\operatorname{Spec}(A)$, lalu tentukan normalisasi
dan dimensi $A$.

\textbf{Petunjuk \oid{O013-M07-H1}.} Unsur $t$ memenuhi polinomial monik
$X^2-t^2\in A[X]$ dan $B=A+At$.

\textbf{Petunjuk \oid{O013-M07-H2}.} Gunakan teorema terletak di atas dan
fakta bahwa perluasan integral mempertahankan dimensi Krull dalam kasus ini.

\textbf{Jawaban \oid{O013-M07-A}.} Persamaan monik menunjukkan $t$ integral;
karena $B=A[t]=A+At$, perluasannya hingga. Teorema terletak di atas membuat
pemetaan spektrum surjektif. Gelanggang $B=k[t]$ tertutup integral dalam
$k(t)$ dan merupakan penutupan integral $A$ di medan pecahannya, sehingga
$B$ adalah normalisasi $A$. Karena $\dim k[t]=1$ dan $A\subset B$ integral,
$\dim A=1$.

\subsection*{\oid{O013-M08}: Nullstellensatz pada parabola afin}

\textbf{Prompt.} Misalkan $k$ tertutup secara aljabar dan
$A=k[x,y]/(y-x^2)$. Klasifikasikan ideal maksimal $A$, jelaskan titik geometri
yang bersesuaian, dan hitung dimensi Krull $A$.

\textbf{Petunjuk \oid{O013-M08-H1}.} Gunakan isomorfisme $A\cong k[x]$ yang
mengirim kelas $y$ ke $x^2$.

\textbf{Petunjuk \oid{O013-M08-H2}.} Terapkan Nullstellensatz pada ideal
maksimal dan tampilkan satu rantai prima dengan panjang satu.

\textbf{Jawaban \oid{O013-M08-A}.} Setiap ideal maksimal berbentuk
\[
\mathfrak m_a=(x-a,y-a^2)/(y-x^2),\qquad a\in k,
\]
dan bersesuaian dengan titik $(a,a^2)$ pada parabola. Isomorfisme dengan
$k[x]$ memberi rantai prima $(0)\subset(x-a)$ dan tidak ada rantai prima yang
lebih panjang. Jadi $\dim A=1$.

\section{Kriteria selesai}

Unit O013 selesai apabila pembelajar:

\begin{enumerate}
\item dapat menggambar ulang DAG dependensi tanpa melihat tabel;
\item memperoleh sedikitnya tujuh dari delapan poin diagnostik setelah
perbaikan terarah;
\item menyelesaikan \oid{O013-M01}--\oid{O013-M03} secara mandiri untuk
jalur R;
\item menyelesaikan sedikitnya empat dari \oid{O013-M04}\allowbreak--\allowbreak\oid{O013-M08},
termasuk satu prompt lokal/modul dan satu prompt spektrum/dimensi, untuk jalur C;
dan
\item dapat menjelaskan dengan contoh mengapa ``representasi sebagai modul''
merupakan jembatan sah, tetapi karakter grup tidak boleh disamakan dengan
titik spektrum gelanggang.
\end{enumerate}

\end{document}
